\histitem{弗雷德霍姆指标与扰动}

指标这一概念于1920年出现在弗里茨·诺特关于积分方程的工作中。在1904年国际数学家大会的演讲中，希尔伯特考察了黎曼在势理论中提出的一个问题：给定平面内一个有界开域 \(\Omega\) ，其边界为光滑简单闭曲线 \(\Gamma\) ，以及三个函数 \(a\) 、\(b\) 、\(f\) ，它们定义在 \(\Gamma\) 上，求一个函数 \(z: \overline{\Omega} \to \mathbf{C}\) ，它在 \(\Omega\) 上全纯、在 \(\overline{\Omega}\) 上连续，并满足

\[
a \mathcal {R} (z) + b \mathcal {I} (z) = f \quad \text { on } \Gamma .
\]

将 \(\Gamma\) 用实数 \(s \in [0, 2\pi]\) 参数化后，希尔伯特证明该问题归结为求解关于 \(\varphi = \mathcal{R}(z)\) 的一个“具有奇异核”的积分方程，即

\[
a (s) \varphi (s) + \int_ {0} ^ {2 \pi} \mathrm{K} (s, t) \varphi (t) d t = f (s).\tag{4}
\]

这里，核 \(\mathrm{K}(s, t)\) 对某个连续函数 A 具有形式 \(b(s) \cot\left(\frac{t-s}{2}\right) + \mathrm{A}(s, t)\) 。因此它沿对角线奇异，上述积分须按柯西主值的意义理解。

诺特 {[}53{]} 指出，与满足弗雷德霍姆择一定理的积分方程不同，对于上述形式的核 K 和非零右端项 \(f\) ，(4) 可能有非平凡解族。他证明解空间由整数

\[
n = \frac {1}{2 \pi} \int_ {\Gamma} d (\log (a - i b)).\tag{5}
\]

若 \(n < 0\) ，则不存在非平凡解；若 \(n \geqslant 0\) ，方程(4)有一个依赖于 2n 个参数的解族。诺特使用“指标”（Index）一词称呼整数 n，并在其中认出一个绕数；在复变函数研究中，这一概念自十九世纪末以来已属经典。

后来的工作继续研究需要扩展弗雷德霍姆理论的方程实例。不过，在弗雷德霍姆映射的概念被系统提炼出来之前，还要再过二十多年；与此同时，对紧算子的扰动的研究也日趋成熟。

1909年，H. 外尔 {[}81{]} 证明，若定义在 \(\ell_{\mathbf{C}}^{2}(\mathbf{N})\) 上的两个有界二次型之差是完全连续的，则它们的本质谱相同（参见 III，第89页，定理2）。此外，里斯在1916年的回忆录中（当然没有使用现代语言）注意到，巴拿赫空间 E 上的紧算子构成 \(\mathcal{L}(\mathrm{E})\) 的双边理想。然而，似乎直到20世纪40年代才有人沿着这两条评注所指示的道路探索。1941年，J. 卡尔金受 J. 冯·诺伊曼关于算子代数的工作启发（我们很快将有机会讨论，见第537页），指出在希尔伯特空间框架下研究模此理想的同余关系的意义 {[}10{]}。他说明了如今以自己名字命名的代数的结构（参见 III，第75页）如何使人能够简单地恢复外尔的结果。

与此同时，大约1942年，J. 迪厄多内 {[}15{]} 对研究赋范空间之间的连续线性映射 \(u\) 的扰动产生兴趣，所考虑的情形是扰动按算子范数的意义“很小”。他只研究 \(u\) 是严格态射且其核有限维或有限余维的情形；他证明 \(u\) 的任意“微小扰动”也具有同样性质，并说明对于弗雷德霍姆映射（参见 III，第40页的 \(cf. n^{\circ} 2\) ），指标局部恒定（III，第58页定理1），III，第55页§4呈现的若干结果也如此。

这两种思想在1950年前后汇合，此时 F. 阿特金森 {[}3{]}、I. 戈赫伯格 {[}26{]}、S. 米赫林 {[}46{]} 和 B. 尤德 {[}93{]} 研究紧算子的扰动。他们明确指出了它与诺特指标的联系，并提出弗雷德霍姆映射 {[}3，第8页{]} 和里斯映射 {[}93，第5节{]} 的概念。随后十年中，人们对后者展开了系统研究；本书第三章呈现的结果在这一时期大致取得了今天的形式。

这些理论通常在巴拿赫空间框架中展开；然而，各种应用促使人们将其推广到更广泛的拓扑向量空间类。因此，弗雷歇空间的情形于1954年自然出现在 H. 嘉当和 J-P. 塞尔的工作中 {[}13{]}，该工作证明紧复解析流形取值于相干层的上同调是有限维的，其证明调用了 L. 施瓦茨证明的变形结果 {[}66{]}（参见 III，第73页定理2）。

在公式(5)的情形，若注意到右端（绕数）在同伦下不变，指标在变形下的不变性便得到充分说明。这样的指标拓扑表达式很快因微分几何中产生的弗雷德霍姆应用而闻名。若 D 是紧 \(C^{\infty}\) 流形 X 上的微分算子（参见 VAR，§14），且 D 是“椭圆的”（参见{[}2，第1节{]}），则将 D 延拓到适当的 Sobolev 空间定义一个弗雷德霍姆映射。寻求一个利用 X 的“特征类”给出其指标的拓扑表达式的公式——这一研究由 I. 盖尔范德约1960年的工作等所提出 {[}23，第I卷，第65页{]}——将在1963年导向 M. 阿蒂亚和 I. 辛格的“指标定理” {[}2{]}。后者将引发分析、拓扑和微分几何交汇处的非凡发展。我们无法在此描述他们开辟的广阔领域；这些领域至今仍极其富有活力，并且不断从弗雷德霍姆理论汲取营养。

